\documentclass[10pt,a4paper]{amsart}
\usepackage[margin=19mm]{geometry}
\usepackage{amsmath,amssymb,mathtools}
\usepackage[scr=rsfs,cal=euler]{mathalfa}
\usepackage{xfrac}
\usepackage[unicode,hidelinks,bookmarks=false]{hyperref}
\newcommand{\ie}{i.e.\ }
\newcommand{\cf}{cf.\ }
\newcommand{\resp}{respectively\ }
\newcommand{\Subsection}{\subsection}
\providecommand{\nobreakdash}{\nobreak\hbox{-}}
\setlength{\emergencystretch}{2em}
\multlinegap0pt
\usepackage{bm}
\usepackage{bbm}
\usepackage{dsfont}
\usepackage{xspace}
\usepackage{cancel}
\usepackage{mathtools}
\usepackage{cleveref}
\usepackage{amsthm}
\makeatletter
\@ifundefined{theorem}{\newtheorem{theorem}{Theorem}}{}
\@ifundefined{proposition}{\newtheorem{proposition}[theorem]{Proposition}}{}
\@ifundefined{remark}{\newtheorem{remark}[theorem]{Remark}}{}
\makeatother
\newcommand{\N}{\mathbb{N}}
\newcommand{\Pell}{\mathcal{P}}
\newcommand{\Z}{\mathbb{Z}}
\title[$SU(2)$-representations of Branched Covers]{$SU(2)$-representations of Branched Covers}
\author{Sudipta Ghosh, Zhenkun Li, Juanita Pinz\'on-Caicedo}
\date{Source: arXiv:2508.19669v1 (CC BY 4.0)}
\begin{document}
\maketitle

\noindent\emph{Source and licence.} $SU(2)$-representations of Branched Covers. Version arXiv:2508.19669v1 (CC BY 4.0), distributed under the Creative Commons Attribution 4.0 International licence, as recorded in the arXiv licence metadata for this exact version and shown on the abstract page https://arxiv.org/abs/2508.19669v1. Full rights notice: licenses/arxiv-2508.19669-CC-BY-4.0.txt.\par\smallskip

\noindent\textit{Editorial scope.} Counted unit: the Proposition whose source label is \texttt{prop: SICUP and Pell's eq} in the article (source0.tex lines 507--513, source label \texttt{prop: SICUP and Pell's eq}), delivered together with the article's own complete original proof of it (source0.tex lines 514--555, one \texttt{proof} environment of 42 lines). No mathematics is written, repaired or re-derived: statement and proof are the article's own text line for line. Required context retained uncounted: prop:circulant (lines 496--503). Every definition, display and label of those lines is retained unchanged. Omitted from this package and not redistributed: 1--495, 504--506, 556--811. Nothing inside a retained excerpt is omitted or altered: every retained body is delivered exactly as the article's lines are written, so no omission receipt of any kind accompanies this package. Citations: the retained excerpts carry no citation command at all, so no bibliography entry of the article is transcribed and no reference is redistributed. Reference adaptation: none was needed and none was made; every reference inside the delivered unit resolves inside it, so no unresolved reference or citation is produced. Printed slips kept literal: no printed word, symbol, hypothesis or number of the counted statement or of its proof was altered. Layout and class: the article's own class is \texttt{amsart} with packages including geometry, graphicx, xcolor, xy, amssymb, enumitem, hyperref, cleveref, MnSymbol, tikz-cd, overpic, caption, todonotes; the wrapper is the lane's pinned \texttt{amsart} editorial wrapper at 10pt on A4, which restates the theorem-like environments the retained text needs and adds the article's own macro definitions that the retained text needs (\texttt{\textbackslash N}, \texttt{\textbackslash Pell}, \texttt{\textbackslash Z}), with extra wrapper packages (bm, bbm, dsfont, xspace, cancel, mathtools, cleveref). The delivered numbering is the unit's own. Assets: no retained span contains a figure environment, a graphics-inclusion command, a PDF-page inclusion, a table or a TikZ picture; no image file is copied into this package, the package declares no asset dependency and no image file is redistributed with it. Licence: version arXiv:2508.19669v1 of this article is distributed under the Creative Commons Attribution 4.0 International licence, as recorded in the arXiv licence metadata for this exact version and shown on the abstract page https://arxiv.org/abs/2508.19669v1. A full rights notice is delivered at licenses/arxiv-2508.19669-CC-BY-4.0.txt. Characters: every retained excerpt is pure ASCII, so the delivered engine, XeLaTeX with its default Latin Modern text font, reports no missing character.
\begin{remark}\label{prop:circulant} Let $C$ be a $d\times d$ SICUP matrix with $d=2r+1$, $r\in\N$. Then we have:
\begin{enumerate}
\item If $\vec{u}$ denotes the vector of ones, then $\vec{u}$ is an eigenvector for $C$ with eigenvalue $\lambda_1=c_1+2(c_2+\ldots c_r)$. Positivity and unimodularity imply that $\lambda_1=1$.
\item There is an eigenvalue $\lambda_j$ of multiplicity 2 for $j=2,\ldots,r$. This allows us to write $\det(C)=\lambda_1(\lambda_2\cdot\ldots\cdot\lambda_r)^2=1$.
%\item $C$ positive definite implies all eigenvalues are positive.
%\item $\lambda_1(\lambda_2\cdot\ldots\cdot\lambda_r)^2=1$
\end{enumerate}
\end{remark}

\begin{proposition}\label{prop: SICUP and Pell's eq}
There is a one-to-one correspondence between the set of $5\times 5$ integral symmetric matrices $\mathcal{M}(5)$ and the set of solutions $(a,b)$ to the generalized Pell's equation $a^{2}-5b^{2}=4$ such that $a\equiv 2~(\text{mod }5)$.
%following two sets
%\[\mathcal{M}(5)\longleftrightarrow \overline{\Pell(5,4)}
%%\{(a,b)\in\mathbb{Z}_{>0}\times\mathbb{Z}~|~a^2=5b^2+4~{\rm and}~a\equiv 2~({\rm mod}~5)\}
%\]
\end{proposition} 

\begin{proof} Notice that any $5\times 5$ integral symmetric circulant matrix can be parametrized by three integers $x$, $l$, and $m$:
\[
\renewcommand{\arraystretch}{1.25}
A(x,l,m) = \left[
\begin{matrix}
	x & l & m & m & l\\
	l & x & l & m & m\\
	m & l & x & l & m\\
	m & m & l & x & l\\
	l & m & m & l & x\\
\end{matrix}
\right].\]
%and its determinant is given by
%\begin{equation*}
%{\rm det}(A(x,l,m)) = (x+2l+2m)(x^2-(l+m)x-l^2+3lm-m^2)^2,
%\end{equation*}
%where, by \Cref{prop:circulant} we have  
A straightforward computation shows that the eigenvalues of $A$ are given by
\begin{align*}
\lambda_1&=x+2l+2m\\
\lambda_2&=\tfrac{1}{2}\left((2x-l-m)+\sqrt{5}(l-m)\right)\\
\lambda_3&=\tfrac{1}{2}\left((2x-l-m)+\sqrt{5}(m-l)\right),
\text{ so that}\\
\det(A)&=\lambda_1\left(\lambda_2\lambda_3\right)^2, \text{ with}\\
\lambda_2\lambda_3&=\tfrac{1}{2}\left((2x-l-m)^2- 5(l-m)^2\right).
\end{align*}
By \Cref{prop:circulant}, ${\rm det}(A)=1$ and $\lambda_1=1$, forcing $\lambda_2\lambda_3= 1$. Thus, for $A$ to be an element of $\mathcal{M}(5)$ we must have 
\begin{equation*}
(2x-l-m)^2- 5(l-m)^2=4.
\end{equation*}
This immediately defines a map
\begin{align*}
\Phi: \mathcal{M}(5)&\longrightarrow\Pell(5,4)=\{(a,b)\in\Z^2 \mid\ a^{2}-5b^{2}=4\}\\
A(x,l,m)&\longrightarrow(2x-l-m,l-m).
\end{align*}
To show that this map is a bijection, first notice that if $(a,b)\in\Pell(5,4)$, then $a\equiv 2 \pmod{5}$. An explicit inverse for $\Phi$ is obtained as 
\begin{align*}
\Phi^{-1}:\Pell(5,4)&\longrightarrow\mathcal{M}(5)\\
(a,b)&\longrightarrow A\left(\frac{2a+1}{5},\frac{2-a+5b}{10},\frac{2-a-5b}{10}\right),
\end{align*}
where we have used the fact that $x=1-2l-2m$.
\end{proof}

\end{document}
